% macros.tex
% Clean macro file for "Coherent presentations for stylic monoids".
%
% Load this file after the packages used by the manuscript, in particular:
% amsmath, amsthm, amssymb, mathtools, fancyhdr, titlesec, hyperref,
% graphicx, and xy with the all, knot, and poly options.

% -----------------------------------------------------------------------------
% Section and theorem styles
% -----------------------------------------------------------------------------

\titleformat{\section}[block]
  {\scshape\filcenter\Large}{\thesection.}{.5em}{}
\titleformat{\subsection}[runin]
  {\bfseries}{\thesubsection.}{.5em}{}[.]
\titleformat{\subsubsection}[runin]
  {\bfseries}{\thesubsubsection.}{.5em}{}[.]
\titlespacing{\subsubsection}{0pt}{10pt}{.5em}

\newtheoremstyle{ntheorem}
  {\topsep}
  {\topsep}
  {\itshape}
  {0pt}
  {\bfseries}
  {.}
  {.5em}
  {\thmname{#1}\thmnumber{ #2}\thmnote{ (#3)}}

\theoremstyle{ntheorem}
\newtheorem{theorem}[subsubsection]{Theorem}
\newtheorem{proposition}[subsubsection]{Proposition}
\newtheorem{lemma}[subsubsection]{Lemma}
\newtheorem{corollary}[subsubsection]{Corollary}
\newtheorem{conjecture}[subsubsection]{Conjecture}

\theoremstyle{definition}
\newtheorem{definition}[subsubsection]{Definition}
\newtheorem{notation}[subsubsection]{Notation}
\newtheorem{example}[subsubsection]{Example}
\newtheorem{remark}[subsubsection]{Remark}
\newtheorem{Algorithm}[subsubsection]{Algorithm}
\newtheorem{property}[subsubsection]{Property}

% -----------------------------------------------------------------------------
% Page layout and running heads
% -----------------------------------------------------------------------------

\pagestyle{fancy}
\setlength{\oddsidemargin}{0cm}
\setlength{\evensidemargin}{0cm}
\setlength{\topmargin}{-1.6cm}
\setlength{\headheight}{1cm}
\setlength{\headsep}{1cm}
\setlength{\textwidth}{16cm}
\setlength{\marginparwidth}{0cm}
\setlength{\footskip}{2cm}
\setlength{\headwidth}{16cm}
\setlength{\arraycolsep}{1pt}

\newcommand{\emptysectionmark}[1]{%
  \markboth{\textbf{#1}}{\textbf{#1}}}
\renewcommand{\sectionmark}[1]{%
  \markboth{\textbf{\thesection. #1}}{\textbf{\thesection. #1}}}
\renewcommand{\subsectionmark}[1]{%
  \markright{\textbf{\thesubsection. #1}}}

\fancyhf{}
\fancyhead[LE]{\leftmark}
\fancyhead[RO]{\rightmark}
\fancyfoot[LE,RO]{\thepage}
\renewcommand{\headrulewidth}{1pt}
\renewcommand{\footrulewidth}{0pt}

\fancypagestyle{plain}{%
  \fancyhf{}%
  \fancyfoot[LE,RO]{\thepage}%
  \renewcommand{\headrulewidth}{0pt}%
  \renewcommand{\footrulewidth}{0pt}%
}

\renewcommand{\headrule}{%
  \hrule width\headwidth height\headrulewidth
  \vskip-\headrulewidth
}

% -----------------------------------------------------------------------------
% Text utilities
% -----------------------------------------------------------------------------

\newcommand{\pdf}[1]{\texorpdfstring{$#1$}{#1}}

\newcount\hh
\newcount\mm
\mm=\time
\hh=\time
\divide\hh by 60
\divide\mm by 60
\multiply\mm by 60
\mm=-\mm
\advance\mm by \time
\newcommand{\hhmm}{\number\hh:\ifnum\mm<10 0\fi\number\mm}

% -----------------------------------------------------------------------------
% Xy-pic setup and the custom four-line arrow variant \ar@4
% -----------------------------------------------------------------------------

\renewcommand{\objectstyle}{\displaystyle}
\renewcommand{\labelstyle}{\displaystyle}
\UseTips
\SelectTips{eu}{11}

\newdir{ >}{{}*!/-10pt/@{>}}
\newdir{ -}{{}*!/-10pt/@{}}
\newdir{> }{{}*!/+10pt/@{>}}

\makeatletter

\xyletcsnamecsname@{dir4{}}{dir{}}
\xydefcsname@{dir4{-}}{\line@ \quadruple@\xydashh@}
\xydefcsname@{dir4{.}}{\point@ \quadruple@\xydashh@}
\xydefcsname@{dir4{~}}{\squiggle@ \quadruple@\xybsqlh@}
\xydefcsname@{dir4{>}}{\Tttip@}
\xydefcsname@{dir4{<}}{\reverseDirection@\Tttip@}

\xydef@\quadruple@#1{%
  \edef\Drop@@{%
    \dimen@=#1\relax
    \dimen@=.5\dimen@
    \A@=-\sinDirection\dimen@
    \B@=\cosDirection\dimen@
    \setboxz@h{%
      \setbox2=\hbox{\kern3\A@\raise3\B@\copy\z@}%
      \dp2=\z@ \ht2=\z@ \wd2=\z@ \box2
      \setbox2=\hbox{\kern\A@\raise\B@\copy\z@}%
      \dp2=\z@ \ht2=\z@ \wd2=\z@ \box2
      \setbox2=\hbox{\kern-\A@\raise-\B@\copy\z@}%
      \dp2=\z@ \ht2=\z@ \wd2=\z@ \box2
      \setbox2=\hbox{\kern-3\A@\raise-3\B@\noexpand\boxz@}%
      \dp2=\z@ \ht2=\z@ \wd2=\z@ \box2
    }%
    \ht\z@=\z@ \dp\z@=\z@ \wd\z@=\z@ \noexpand\styledboxz@
  }%
}

\xydef@\Tttip@{%
  \kern2pt \vrule height2pt depth2pt width\z@
  \Tttip@@ \kern2pt \egroup
  \U@c=0pt \D@c=0pt \L@c=0pt \R@c=0pt \Edge@c={\circleEdge}%
  \def\Leftness@{.5}\def\Upness@{.5}%
  \def\Drop@@{\styledboxz@}%
  \def\Connect@@{\straight@{\dottedSpread@\jot}}%
}

\xydef@\Tttip@@{%
  \dimen@=.25\dimen@
  \B@=\cosDirection\dimen@
  \setboxz@h\bgroup
  \reverseDirection@\line@
  \wdz@=\z@ \ht\z@=\z@ \dp\z@=\z@
  {\vDirection@(1,-1)\xydashl@ \xyatipfont\char\DirectionChar}%
  {\vDirection@(1,+1)\xydashl@ \xybtipfont\char\DirectionChar}%
}

% Extend Xy-pic's arrow parser with variant 4.  This is deliberately a direct
% definition: \xydef@ would emit a warning because \ar@form already exists.
\gdef\ar@form{%
  \ifx \space@\next \expandafter\DN@\space{\xyFN@\ar@form}%
  \else\ifx ^\next \DN@ ^{\xyFN@\ar@style}\edef\arvariant@@{\string^}%
  \else\ifx _\next \DN@ _{\xyFN@\ar@style}\edef\arvariant@@{\string_}%
  \else\ifx 0\next \DN@ 0{\xyFN@\ar@style}\def\arvariant@@{0}%
  \else\ifx 1\next \DN@ 1{\xyFN@\ar@style}\def\arvariant@@{1}%
  \else\ifx 2\next \DN@ 2{\xyFN@\ar@style}\def\arvariant@@{2}%
  \else\ifx 3\next \DN@ 3{\xyFN@\ar@style}\def\arvariant@@{3}%
  \else\ifx 4\next \DN@ 4{\xyFN@\ar@style}\def\arvariant@@{4}%
  \else\ifx \bgroup\next \let\next@=\ar@style
  \else\ifx [\next \DN@[##1]{\ar@modifiers{[##1]}}%]
  \else\ifx *\next \DN@ *{\ar@modifiers}%
  \else\addLT@\ifx\next \let\next@=\ar@slide
  \else\ifx /\next \let\next@=\ar@curveslash
  \else\ifx (\next \let\next@=\ar@curveinout %)
  \else\addRQ@\ifx\next \addRQ@\DN@{\ar@curve@}%
  \else\addLQ@\ifx\next \addLQ@\DN@{\xyFN@\ar@curve}%
  \else\addDASH@\ifx\next \addDASH@\DN@{\defarstem@-\xyFN@\ar@}%
  \else\addEQ@\ifx\next \addEQ@\DN@{\def\arvariant@@{2}\defarstem@-\xyFN@\ar@}%
  \else\addDOT@\ifx\next \addDOT@\DN@{\defarstem@.\xyFN@\ar@}%
  \else\ifx :\next \DN@:{\def\arvariant@@{2}\defarstem@.\xyFN@\ar@}%
  \else\ifx ~\next \DN@~{\defarstem@~\xyFN@\ar@}%
  \else\ifx !\next \DN@!{\dasharstem@\xyFN@\ar@}%
  \else\ifx ?\next \DN@?{\ar@upsidedown\xyFN@\ar@}%
  \else \let\next@=\ar@error
  \fi\fi\fi\fi\fi\fi\fi\fi\fi\fi\fi\fi
  \fi\fi\fi\fi\fi\fi\fi\fi\fi\fi\fi \next@
}

\makeatother

% -----------------------------------------------------------------------------
% Arrows and delimiters
% -----------------------------------------------------------------------------

\newcommand{\fl}{\rightarrow}
\newcommand{\flg}{\leftarrow}
\newcommand{\fll}{\longrightarrow}
\newcommand{\ofl}[1]{\overset{\displaystyle #1}{\fll}}
\newcommand{\ifl}{\rightarrowtail}
\newcommand{\pfl}{\twoheadrightarrow}
\newcommand{\fllg}{\longleftarrow}
\newcommand{\oflg}[1]{\overset{\displaystyle #1}{\fllg}}
\newcommand{\dfl}{\Rightarrow}
\newcommand{\dfll}{\Longrightarrow}
\newcommand{\odfl}[1]{\overset{\displaystyle #1}{\dfll}}
\newcommand{\tfl}{\Rrightarrow}
\newcommand{\qfl}{\xymatrix@1@C=10pt{\ar@4[r]&}}

\newcommand{\ens}[1]{\left\{#1\right\}}
\renewcommand{\angle}[1]{\left\langle#1\right\rangle}
\newcommand{\pres}[2]{\langle #1\mid #2\rangle}
\newcommand{\bigpres}[2]{\bigl\langle #1\bigm|#2\bigr\rangle}
\newcommand{\cohpres}[3]{\langle #1\mid #2\mid #3\rangle}
\newcommand{\bigcohpres}[3]{%
  \bigl\langle #1\bigm|#2\bigm|#3\bigr\rangle}

% -----------------------------------------------------------------------------
% Accents, operators, and notation
% -----------------------------------------------------------------------------

\newcommand{\op}[1]{#1^{\mathrm{o}}}
\newcommand{\env}[1]{#1^{\mathrm{e}}}
\newcommand{\cl}[1]{\overline{#1}}
\newcommand{\rep}[1]{\widehat{#1}}
\renewcommand{\tilde}[1]{\widetilde{#1}}
\newcommand{\tck}[1]{#1^{\top}}
\newcommand{\vect}[1]{\overrightarrow{#1}}

\DeclareMathOperator{\id}{Id}
\DeclareMathOperator{\Ker}{Ker}
\DeclareMathOperator{\coker}{coker}
\DeclareMathOperator{\Coker}{Coker}
\DeclareMathOperator{\End}{End}
\DeclareMathOperator{\Aut}{Aut}
\DeclareMathOperator{\Lan}{Lan}
\DeclareMathOperator{\h}{h}
\DeclareMathOperator{\crit}{c}
\DeclareMathOperator{\Hom}{Hom}
\DeclareMathOperator{\Ext}{Ext}
\DeclareMathOperator{\Tor}{Tor}
\DeclareMathOperator{\Der}{Der}
\DeclareMathOperator{\FDT}{FDT}
\DeclareMathOperator{\FDTAB}{FDT_{\mathrm{ab}}}
\DeclareMathOperator{\FP}{FP}
\DeclareMathOperator{\FHT}{FHT}
\DeclareMathOperator{\U}{\mathcal{U}}
\DeclareMathOperator{\Tc}{\mathcal{T}}
\DeclareMathOperator{\Si}{\mathcal{S}}
\DeclareMathOperator{\Grw}{Grw}
\DeclareMathOperator{\Rec}{Rect}
\DeclareMathOperator{\Sh}{S}
\DeclareMathOperator{\ShST}{ShSt}
\DeclareMathOperator{\ShT}{ShT}
\DeclareMathOperator{\YoungST}{Ys}
\DeclareMathOperator{\YoungT}{Yt}
\DeclareMathOperator{\SkewT}{St}
\DeclareMathOperator{\RSK}{sRSK}
\DeclareMathOperator{\PreRow}{PreRow}
\DeclareMathOperator{\col}{Col}
\DeclareMathOperator{\Ba}{Ba}
\DeclareMathOperator{\ir}{ir}
\DeclareMathOperator{\irm}{ir_{\mathrm{m}}}
\DeclareMathOperator{\lm}{lm}
\DeclareMathOperator{\lc}{lc}
\DeclareMathOperator{\lt}{lt}
\DeclareMathOperator{\wt}{wt}
\DeclareMathOperator{\up}{up}
\DeclareMathOperator{\down}{down}
\DeclareMathOperator{\sh}{sh}
\DeclareMathOperator{\std}{std}
\DeclareMathOperator{\height}{ht}
\DeclareMathOperator{\rowword}{r}
\DeclareMathOperator{\NF}{NF}
\DeclareMathOperator{\Crit}{Crit}
\DeclareMathOperator{\rem}{Rem}

\newcommand{\styl}{\mathrm{styl}}
\newcommand{\plac}{\mathrm{plac}}
\newcommand{\supp}{\operatorname{supp}}
\newcommand{\Styl}{\mathrm{Styl}}
\newcommand{\NTab}{\mathrm{NTab}}
\newcommand{\lex}{\mathrm{lex}}
\newcommand{\slex}{\mathrm{slex}}
\newcommand{\Irr}{\mathrm{Irr}}
\newcommand{\deglex}{\mathrm{deglex}}
\newcommand{\RRightarrow}{%
  \mathrel{\vcenter{\hbox{\ensuremath{\Rightarrow}}}}}
\newcommand{\tens}{\otimes}
\newcommand{\Nb}{\mathbb{N}}
\newcommand{\Zb}{\mathbb{Z}}
\newcommand{\dr}{\partial}

\renewcommand{\phi}{\varphi}
\renewcommand{\epsilon}{\varepsilon}
\renewcommand{\leq}{\leqslant}
\renewcommand{\geq}{\geqslant}

% Calligraphic symbols.  These definitions use \mathcal consistently and do
% not require the separate euscript package.
\newcommand{\Ar}{\mathcal{A}}
\newcommand{\Br}{\mathcal{B}}
\newcommand{\Cr}{\mathcal{C}}
\newcommand{\Dr}{\mathcal{D}}
\newcommand{\Er}{\mathcal{E}}
\newcommand{\Fr}{\mathcal{F}}
\newcommand{\Gr}{\mathcal{G}}
\newcommand{\Hr}{H_r}
\newcommand{\Hl}{H_l}
\newcommand{\Ir}{\mathcal{I}}
\newcommand{\Kr}{\mathcal{K}}
\newcommand{\Lr}{\mathcal{L}}
\newcommand{\Mr}{\mathcal{M}}
\newcommand{\Nr}{\mathcal{N}}
\newcommand{\Or}{\mathcal{O}}
\renewcommand{\Pr}{\mathcal{P}}
\newcommand{\Qr}{\mathcal{Q}}
\newcommand{\Rr}{\mathcal{R}}
\newcommand{\Sr}{\mathcal{S}}
\newcommand{\Tr}{\mathcal{T}}
\newcommand{\Ur}{\mathcal{U}}
\newcommand{\Vr}{\mathcal{V}}
\newcommand{\Yr}{\mathcal{Y}}
\newcommand{\Xr}{\mathcal{X}}

% Bold symbols.
\newcommand{\A}{\mathbf{A}}
\newcommand{\B}{\mathbf{B}}
\newcommand{\C}{\mathbf{C}}
\newcommand{\D}{\mathbf{D}}
\newcommand{\F}{\mathbf{F}}
\newcommand{\G}{\mathbf{G}}
\newcommand{\K}{\mathbb{K}}
\newcommand{\M}{\mathbf{M}}
\renewcommand{\S}{\mathbf{S}}
\newcommand{\V}{\mathbf{V}}
\newcommand{\W}{\mathbf{W}}
\newcommand{\ii}{\mathbf{i}}
\providecommand{\R}{}
\renewcommand{\R}{\mathbf{R}}
\renewcommand{\P}{\mathbf{P}}
\providecommand{\T}{}
\renewcommand{\T}{\mathbf{T}}
\providecommand{\Q}{}
\renewcommand{\Q}{\mathbf{Q}}
\newcommand{\Sk}{\mathbf{Sk}}
\newcommand{\St}{\mathbf{Std}}

% Categories.
\newcommand{\catego}[1]{\mathbf{\mathsf{#1}}}
\newcommand{\Ab}{\catego{Ab}}
\newcommand{\Gp}{\mathbf{Gp}}
\newcommand{\Cat}{\catego{Cat}}
\newcommand{\Act}{\catego{Act}}
\newcommand{\Mod}{\catego{Mod}}
\newcommand{\Nat}{\catego{Nat}}
\newcommand{\Set}{\catego{Set}}
\newcommand{\Ord}{\mathbf{Ord}}
\newcommand{\Rep}{\catego{Rep}}
\newcommand{\LinCat}{\catego{LinCat}}
\newcommand{\Vect}{\catego{Vect}}
\newcommand{\Vectg}{\catego{GrVect}}
\newcommand{\Alg}{\catego{Alg}}
\newcommand{\Algg}{\catego{GrAlg}}
\newcommand{\Ch}{\catego{Ch}}
\newcommand{\Chg}{\catego{GrCh}}
\newcommand{\Grph}{\catego{Grph}}

\newcommand{\Ct}[1]{\mathrm{Ct}(#1)}
\newcommand{\Wk}[1]{\mathrm{Wk}(#1)}
\newcommand{\As}{\mathbf{As}}
\newcommand{\Mon}{\mathbf{Mon}}
\newcommand{\Perm}{\mathbf{Perm}}
\newcommand{\Brd}{\mathbf{Braid}}
\newcommand{\Colo}{\mathbf{Coln}}
\newcommand{\Std}{\mathbf{Std}}

% Cohomological notation.
\newcommand{\fnat}[2]{F_{#1}[#2]}
\newcommand{\ho}{\mathrm{H}}
\newcommand{\poldim}[1]{d_{\mathrm{pol}}(#1)}
\newcommand{\cohdim}[1]{\mathrm{cd}(#1)}

% Miscellaneous notation used in rewriting diagrams.
\renewcommand{\labelitemi}{$-$}
\renewcommand{\theenumi}{\roman{enumi}}
\renewcommand{\labelenumi}{\textup{\textbf{\theenumi)}}}
\newcommand{\tableau}[1]{\texttt{#1}}
\newcommand{\Cdot}{\!\cdot\!}
\newcommand{\Pl}{\mathbf{P}}
\newcommand{\HP}{\mathbf{Hp}}
\newcommand{\insl}[1]{\rightsquigarrow_{#1}}
\newcommand{\insr}[1]{%
  \;\raisebox{0.1em}{\rotatebox[origin=c]{180}{$\rightsquigarrow$}}_{#1}\;}
\newcommand{\Right}[1]{\ensuremath{\xrightarrow{\,#1\,}}}
\newcommand{\Left}[1]{\ensuremath{\xleftarrow{\,#1\,}}}
\newcommand{\Up}[1]{%
  \ensuremath{\big\uparrow\!\text{\raisebox{.1ex}{\scriptsize $#1$}}}}
\newcommand{\Down}[1]{%
  \ensuremath{\big\downarrow\!\text{\raisebox{.1ex}{\scriptsize $#1$}}}}
\newcommand{\DRight}[1]{\ensuremath{\xdashrightarrow{\,#1\,}}}
\newcommand{\len}{\ell}
\newcommand{\Ich}{C_r}
\newcommand{\Jch}{C_l}
\newcommand{\points}{\rotatebox{90}{\hspace{-1pt}\ldots}}

\makeatletter
\newcommand{\blfootnote}{\xdef\@thefnmark{}\@footnotetext}
\makeatother
